\documentclass[a4paper,twocolumn]{article}
\usepackage[T1]{fontenc}
\usepackage[utf8]{inputenc}
\usepackage[dvips]{graphicx}
\usepackage[dutch]{babel}
\begin{document}

\title{Het Semmelweis effect}

\author{Roel Pompen}
\maketitle

% \section*{\includegraphics{./TODO/boek/roel.eps}}
\includegraphics{./TODO/boek/roel.eps}
\begin{abstract}
Recente contradicties in getallen die betrekking hebben op de COVID-19 pandemie dreigen de crisis een ander karakter te geven.
Om het mogelijk te maken deze nieuwe gegevens, en
de situatie waarin we ons \"uberhaupt bevinden, te interpreteren duik
ik even de medische geschiedenis in. Op basis van dit verhaal hoop
ik dat mensen gerustgesteld worden als ze geconfronteerd worden met
cijfers die simpelweg niet compleet zijn en derhalve geen interpretatie
toelaten. Niet de fijnste vorm van geruststelling, maar we zullen
het er in deze tijden mee moeten doen.
\end{abstract}

\section*{Hygi\"ene in het ziekenhuis}

In 1847 in het algemeen ziekenhuis van Wenen, waren er twee kraamafdelingen:
E\'en gerund door doctoren en eentje door verpleegsters. Bij de doktoren
was het sterftecijfer drie keer zo hoog als bij de vleugel waar verpleegsters
de scepter zwaaiden en dit al meer dan een eeuw. (Het ziekenhuis stamt
uit ca. 1700)

In een tijd van 2 maanden zakte het sterftecijfer bij de door doktoren
beheerde vleugel tot op een lager niveau dan bij de verpleegsters.
Na een jaar zelfs tot bijna nul. Dit fenomenale resultaat was het
gevolg van de acties van Ignaz Semmelweis die hygi\"enische maatregelen
invoerde. En dit nog voordat beroemdheden als Louis Pasteur en Joseph
Lister (waarnaar het Listerine mondwater is vernoemd), zich bewust
werden van microben en derhalve een verklaring konden verbinden aan
de effectiviteit van deze strategie.

\subsection*{De cijfers}

Op de door doktoren beheerde vleugel stierven jaarlijks 10\% van de
baby's en moeders. Bij de door verpleegsters beheerde vleugel was
dit 4\%. De reden was dat chirurgen die in de kelder les gaven, gebruik
makend van kadavers die vaak vol infecties zaten, zonder hun handen
te wassen en vergezeld van hun vuile instrumenten, later op de dag
assistentie gingen verlenen op de kraamafdeling.

Ignaz Semmelweis stelde voor dat instrumenten voor autopsie en anatomische
les in de kelder moesten blijven en dat men, bij verplaatsing tussen
de vertrekken, de handen moest wassen in chloorkalk. Ook de instrumenten
moesten gereinigd worden. Het gevolg was dat na 14 maanden het sterftecijfer
op de kraamafdeling tot praktisch nul was gezakt.

\subsection*{Academische belediging}

Semmelweis werd niet als held gezien, omdat men het ziekenhuis met
terugwerkende kracht ging verwijten door gebrek aan hygi\"ene verantwoordelijk
te zijn voor vele doden. Iedereen had in die tijd nabije familieleden
verloren in het kraambed. Dus werd Semmelweis verketterd. Hij stierf
doodongelukkig en alleen. Zijn levensloop is niet precies bekend.
Hij belandde in het gesticht, maar de opgaaf van reden wordt in het
algemeen betwist. Werd hij gek ten gevolge van syfilis? Werd hij depressief
van de vijandigheid tegen hem? Feit blijft dat hij in de krochten
van een 19\textsuperscript{e}-eeuws gesticht stierf en zijn tijdgenoten
hem kapot schreven. Hij stierf overigens ten gevolge van gangreen
(weefselsterfte), na een afranseling door bewakers.

\section*{Het Semmelweis effect}

Nu is de titel van dit artikel verduidelijkt: Het Semmelweis effect
noem ik het verschijnsel dat als er middels een simpele verbetering
een grote vooruitgang wordt geboekt, men de bedenkers ervan verkettert,
\'of dat de vooruitgang wordt tegengewerkt tot de gevestigde orde met
pensioen is, of is overleden.

Ik ga proberen aan te tonen dat de gevestigde orde handelt uit commercieel
belang en wetenschappelijke dwaling en dat er een tijd gaat aanbreken
waarin de huidige generatie ``experts''
uit het mediabeeld mag vertrekken, hopelijk nog v\'o\'or hun pensioen;
Een positief gevolg van de corona crisis.

Overigens doet dit Semmelweis effect zich in meerdere wetenschappelijke
disciplines voor, zelfs daar waar je het niet zou verwachten; In exacte
wetenschappen zoals Wiskunde en Informatica. De oververtegenwoordiging
van de medische sector is het gevolg van het feit dat artsen geremd
zijn in dingen in twijfel te trekken; Men leert immers dat als je
op school niet goed oplet en niet alles doet zoals het geleerd is
er doden kunnen vallen. Daar tegenover staan bijvoorbeeld theoretische
natuurkundigen die bereid moeten zijn om alles wat ze geleerd hebben
in twijfel te trekken, anders kunnen ze hun vak zelfs niet goed beoefenen!

\section*{De relatie met COVID-19}

Men begint langzaam te begrijpen dat we geadviseerd willen worden
door experts die g\'e\'en banden hebben met de commerci\"ele sector. Een
epidemioloog heeft in de vrije markt niets te zoeken. We weten vanuit
de epidemiologie dat een lockdown niet kan werken, hetgeen ook blijkt
uit de laatste cijfers. Ook is men nu verbaasd dat er meer resistente
mensen zijn dan verwacht. Zoals eerder gesteld, ga ik in dit artikel
geen cijfers noemen. Die veranderen immers continue. W\'el maak ik inschattingen
en doe ik beweringen die binnen nu en twee jaar bevestigd worden of
worden ontkracht. Ik schrijf dit artikel natuurlijk omdat ik verwacht
dat mijn cijfers en beweringen in de nu volgende paragraaf overeind
blijven staan.

\subsection*{Epidemiologische aspecten}

Gebieden zonder lockdown en met lockdown doen niet voor elkaar onder
en dit was absoluut te verwachten. Uitgaande van een lockdown effectiviteit
van 95\% (een veel te ruime schatting) en een basaal reproductief
getal van 2 en een generatie interval van 5,2 dagen (volgens de huidige
metingen), vertragen we de epidemie met 22 dagen.

Daar staat tegenover
dat je niet weet waar de nieuwe brandhaarden zich topografisch bevinden,
met andere woorden de prijs die je voor de vertraging betaalt is dat
je pas na een dag of 10 in kaart hebt waar de epidemie zich continueert.
We vechten dan dus bijna twee weken blind. En we zitten dan dus langer
in lockdown. 

\subsection*{Averechts}

De draconische maatregelen die worden getroffen op initiatief van
burgers en bedrijven werken vaak averechts, omdat deze strategie\"en
(zonder medische opleiding) niet correct worden uitgevoerd. Men kan
van de gewone burger niet verwachten dat maatregelen zoals handschoenen
correct worden toegepast.

En handschoenen zijn een potentieel gevaar. De huid is onze eerste
barri\`ere tegen infecties en maakt een groep stoffen aan, defensines,
melkzuur en vetzuren geheten, die microben neutraliseren. Van de andere kant bekeken is het echter ook zo
dat de temperatuur van onze huid de vermenigvuldiging van
microben doet versnellen.

Door handschoenen te dragen haal je de neutraliserende
effecten weg en houdt je de lichaamswarmte over. Je kunt nu op je
handschoenen microben opsparen en die uitsmeren in je bloedbaan, zo
dicht mogelijk bij de hersenen, want: Een mens zit per dag gemiddeld
320 keer aan zijn gezicht, waar slijmvliezen bij de mond, ogen en
neus de toegang vormen naar de bloedbaan.

Mocht je geen handschoenen dragen maar je handen wel 5 keer per dag
reinigen met desinfecterende gel, dan komt defensineproductie in gevaar
en is de huid op je handen zijn natuurlijke barri\`ere tegen bacteri\"en
kwijt.

\subsection*{Schadelijk burgerinitiatief}
Het feit dat huismeesters in flatgebouwen de liften deactiveren zodat
iedereen met zijn handen dezelfde trapleuning moet aanraken helpt
ook niet echt (tenminste 50\% van de mensen boven de 45 kan niet met
een zware tas boodschappen de trap op zonder de leuning te gebruiken).
Bij supermarkten wordt je verplicht een karretje uit een collectie
van 80 te pakken, want allemaal gedwongen dezelfde objecten moeten
aanraken is natuurlijk een goed idee\ldots{} En zo kan ik nog wel
even doorgaan.

Hoewel men zal zeggen dat met structurele ontsmetting wordt voorkomen
dat er iets misgaat, is dat andersom geredeneerd; Raak niet allemaal
dezelfde paar karretjes aan, hopende dat deze met vlijt en discipline
ontsmet worden. De ontsmetting lost immers een probleem op dat je,
als je slechts een brood komt kopen, helemaal niet hebt!

\subsection*{Superinfecties}

De reden dat dit verhaal over bacteri\"ele infecties relevant is, is
dat mensen met m\'e\'er dan een serieuze infectie (een superinfectie)
binnen de risicogroep zitten, met andere woorden: Als je een bacteri\"ele
infectie hebt en je krijgt er een corona infectie bij (of andersom)
dan krijgt je immuunsysteem meer moeite om je in leven te houden.

\subsection*{Commercieel belang}

We weten sinds het werk van de grootvader van de epidemiologie (William
Farr, belangrijkste publicaties rond 1870) dat een epidemie met een
exponenti\"ele stijging begint en afvlakt als verspreiding niet meer
lukt. Dit is op het moment dat er teveel mensen reeds resistent zijn. 

Virologen proberen met nieuwe technologie (medicijnen en testen) deze
curve te be\"invloeden, maar wiskundigen weten dat dat maar beperkt
kan lukken. Investeerders snappen dat echter niet en voil\'a, een markt
is geboren.

Het zijn d\'eze experts die de gevestigde orde vormen die onderhand
het veld moet gaan ruimen. We moeten ons laten adviseren door onafhankelijke
specialisten die zich baseren op ongecompromitteerde wetenschap. Niet
de wetenschap van de plagiaat generatie. Die is immers voor een groot
deel gebaseerd op resultaatgericht onderzoek. 

\subsection*{Opportunisme}

In deze tijden struikelen de mensen van de farmaceutische industrie
over elkaar om funding binnen te harken voor hun medicijn, dat dan
over 2 jaar, dus n\'a de epidemie beschikbaar wordt en \emph{misschien
}voor de volgende mutatie \'o\'ok wel werkt.

Zelfs in niet-medische takken van sport zien we een idiote hype. In
Nederland ontwikkelt men een corona-app zonder de specificaties duidelijk
te hebben, te weten of er behoefte aan is en of het in de praktijk
wel kan werken. Zolang er iemand de financiering voor regelt is alles
goed.

\subsection*{Verkeerd advies en cijfers}

\begin{table}
\noindent \begin{centering}
\begin{tabular}{|c|c|c|}
\hline 
\multicolumn{3}{|c|}{aantal / \% bevolking getest}\tabularnewline
\hline 
\textbf{seronegatief} &  & \textbf{seropositief}\tabularnewline
\hline 
aantal / \% &  & aantal / \%\tabularnewline
\hline 
\multicolumn{3}{|c|}{\textbf{buiten risicogroep}}\tabularnewline
\hline 
\textbf{hospitalisaties} & \textbf{overleden} & \textbf{genezen}\tabularnewline
\hline 
aantal / \% & aantal / \% & aantal / \%\tabularnewline
\hline 
\multicolumn{3}{|c|}{\textbf{binnen risicogroep}}\tabularnewline
\hline 
\textbf{hospitalisaties} & \textbf{overleden} & \textbf{genezen}\tabularnewline
\hline 
aantal / \% & aantal / \% & aantal / \%\tabularnewline
\hline 
\end{tabular}
\par\end{centering}
\caption{Momentopname van een epidemie}
\end{table}

Een epidemioloog die in de media commentaar geeft op cijfers heeft
de academische en morele verplichting de cijfers af te keuren omdat
ze incompleet zijn en dient te zorgen dat hij of zij slechts aan de
hand van het volledige plaatje commentaar geeft. Een suggestie voor
het omschrijven van de toestand op enig moment heb ik in een tabel
bijgevoegd.

Natuurlijk is er naast zo'n tabel ook een omschrijving van de risicogroep
nodig om te voorkomen dat er overledenen worden opgeteld bij de infectieziekte,
terwijl het eigenlijk net zo goed aan een onderliggende conditie te
verwijten valt. Zo is het aantal overlijdens door hartaanvallen gedaald
ten opzichte van normaal, omdat deze mensen bij de corona slachtoffers
worden geteld!

\section*{Conclusie}

Een expert die complete presentatie van feiten niet ondersteunt cre\"eert
onzekerheid. De onzekerheid leidt tot paniek. De paniek leidt tot
druk op de politiek. De politiek gaat vervolgens afwijken van de normaal
geachte epidemiologische strategie (``draconische maatregelen'').
De burger neemt initiatieven die aansluiten bij de intensiviteit van
de reactie van de overheid. En voil\'a\ldots

Het is niet zo dat men massaal naar het ziekenhuis komt om te sterven.
Er sterven meer mensen omdat men massaal naar het ziekenhuis komt\ldots
\'of andere plaatsen. Bange mensen gaan clusteren. Bij de apotheek,
want de lockdown veroorzaakt kwaaltjes. Bij de kerk (In Nederland
mogen kerkelijke samenscholingen tot 30 personen w\'el). Lockdown parties.
Parken en bossen.

Wat is er gaande? We weten niet meer wie de experts zijn. Overheden
moeten overheids-epidemiologen aanstellen die de strategie bepalen
en verplicht via zowel commerci\"ele als publieke media minstens twee
keer per week met complete cijfers verslag uitbrengen. Zodoende nemen
ze de ons welbekende experts uit de media, met hun opmerkelijke omissies,
de wind uit de zeilen. Dit doorbreekt de genoemde vicieuze cirkel
van paniek en politiek bijsturen, zodat we kunnen voorkomen dat zoiets
nog eens gebeurt.
\end{document}
